\subsection{Complex Linear Representations}

In this subsection, we assume that K is the field C of complex numbers.

Let \((M,\pi)\) be a linear representation of G. We say that a Hermitian form \(\Phi\) on M is invariant under G if we have

\[
\Phi (\pi (g) x, \pi (g) x ^ {\prime}) = \Phi (x, x ^ {\prime})\tag{48}
\]

for all \(x, x' \in M\) and every \(g \in G\) . This also means that for every \(g \in G\) , the automorphism \(\pi(g)\) of M is unitary with respect to \(\Phi\) .

PROPOSITION 11. --- Let \((M,\pi)\) be a finite-dimensional linear representation of G.

\begin{enumerate}
\def\labelenumi{\alph{enumi})}
\item
  There exists on M a Hermitian form that is positive, separating, and invariant under G.
\item
  Suppose that the representation \(\pi\) is simple. If \(\Phi\) and \(\Psi\) are nonzero Hermitian forms on M that are invariant under G, then there exists a real number a such that \(\Psi = a\Phi\) .
\end{enumerate}

Choose a separating positive Hermitian form on the vector space M, and denote it by \(\Phi_{0}\) . We define a separating positive Hermitian form \(\Phi\) that is invariant under G by setting

\[
\Phi (x, x ^ {\prime}) = \sum_ {g \in \mathrm{G}} \Phi_ {0} (\pi (g) x, \pi (g) x ^ {\prime})\tag{49}
\]

for \(x, x' \in M\) .

Let \(\Psi\) be a Hermitian form on M; there exists a unique endomorphism A of M such that \(\Psi(x, x') = \Phi(x, Ax')\) for \(x, x'\) in M. If, moreover, \(\Psi\) is invariant under G, then the endomorphism A commutes with the automorphism \(\pi(g)\) for \(g \in G\) . If the representation \(\pi\) is simple, then by Schur's lemma (VIII, p.~47, Theorem 1), A is a homothety and there consequently exists a complex number \(a\) such that \(\Psi = a\Phi\) . Since \(\Phi\) and \(\Psi\) are Hermitian and \(\Phi\) is nonzero, \(a\) is a real number. The proposition follows.

We endow the vector space \(\mathbf{C}[\mathrm{G}]\) of complex functions on \(\mathrm{G}\) with the Hilbert space structure whose inner product is given by

\[
\langle f | f ^ {\prime} \rangle_ {\mathrm{G}} = | \mathrm{G} | ^ {- 1} \sum_ {g \in \mathrm{G}} \overline {{f (g)}} f ^ {\prime} (g).\tag{50}
\]

For any function \(f \in C[G]\) , we denote by \(f^{*}\) the function defined by

\[
f ^ {*} (g) = \overline {{f (g ^ {- 1})}}\tag{51}
\]

for \(g \in \mathbf{G}\) ; the mapping \(f \mapsto f^*\) is a semilinear involution of \(\mathbf{C}[\mathbf{G}]\) . We also have

\[
\langle f | f ^ {\prime} \rangle_ {\mathrm{G}} = \langle f ^ {*}, f ^ {\prime} \rangle_ {\mathrm{G}}\tag{52}
\]

for \(f, f' \in \mathbf{C}[\mathrm{G}]\) , with the notation of formula (28) of VIII, p.~410. We therefore have

\[
\langle f | f ^ {\prime} \rangle_ {\mathrm{G}} = | \mathrm{G} | ^ {- 2} \tau (f ^ {*} f ^ {\prime}).\tag{53}
\]

Let \((\mathrm{M},\pi)\) be a finite-dimensional linear representation of G. We endow the vector space M with the structure of a Hilbert space for which the endomorphisms \(\pi(g)\) are unitary (Proposition 11). If we denote by \(A^{*}\) the adjoint of a endomorphism A of M for this structure, then we have \(\mathrm{Tr}(A^{*})=\overline{\mathrm{Tr}(A)}\) . For every \(g\in G\) , we have \(\pi(g^{-1})=\pi(g)^{*}\) , and therefore \(\chi_{\pi}(g^{-1})=\overline{\chi_{\pi}(g)}\) ; in other words, we have \(\chi_{\pi}=\chi_{\pi}^{*}\) . The orthogonality relation for characters (VIII, p.~410, Proposition 4) then has the form

\[
\langle \chi_ {\lambda} | \chi_ {\mu} \rangle_ {\mathrm{G}} = \delta_ {\lambda \mu}\tag{54}
\]

for \(\lambda, \mu \in \widehat{\mathbf{G}}\) . It expresses the fact that the family of the characters \((\chi_{\lambda})_{\lambda \in \widehat{\mathbf{G}}}\) of the simple representations of \(\mathbf{G}\) is an orthonormal basis of the Hilbert space \(\mathrm{Z}(\mathbf{C}[\mathrm{G}])\) of central functions.

Let \(\pi\) and \(\pi'\) be finite-dimensional linear representations of G. We have the relation \(\langle \chi_{\pi}|\chi_{\pi'}\rangle_{\mathrm{G}} = \dim_{\mathbf{C}}\mathrm{Hom}_{\mathrm{G}}(\pi, \pi')\) (VIII, p.~410, Corollary). The representation \(\pi\) is irreducible if and only if \(\langle \chi_{\pi}|\chi_{\pi}\rangle_{\mathrm{G}} = 1\) .

For every element \(\lambda\) of \(\widehat{G}\) , we endow the vector space \(V_{\lambda}\) with the structure of a Hilbert space for which the automorphisms \(\pi_{\lambda}(g)\) are unitary. We denote by \(\langle v|v'\rangle_{\lambda}\) the inner product of two elements \(v, v'\) of \(V_{\lambda}\) and by \(u^{*}\) the adjoint of an endomorphism \(u\) of \(V_{\lambda}\) . Let \(A = (A_{\lambda})_{\lambda \in \widehat{G}}\) and \(A' = (A_{\lambda}')_{\lambda \in \widehat{G}}\) be elements of \(\mathrm{F}(\widehat{\mathrm{G}})\) . We write \(\mathrm{A}^* = (\mathrm{A}_\lambda^*)_{\lambda \in \widehat{\mathrm{G}}}\) . We have \(\overline{\mathcal{F}} (a^{*}) = (\overline{\mathcal{F}} (a))^*\) for every element \(a\) of \(\mathbf{C}[\mathrm{G}]\) . Set

\[
\langle \mathrm{A} | \mathrm{A} ^ {\prime} \rangle_ {\widehat {\mathrm{G}}} = | \mathrm{G} | ^ {- 2} \widehat {\tau} (\mathrm{A} ^ {*} \mathrm{A} ^ {\prime}).\tag{55}
\]

By formula (10) of VIII, p.~407, we have

\[
\langle \mathrm{A} | \mathrm{A} ^ {\prime} \rangle_ {\widehat {\mathrm{G}}} = \frac {1}{| \mathrm{G} | ^ {2}} \sum_ {\lambda \in \widehat {\mathrm{G}}} d _ {\lambda} \operatorname{Tr} (\mathrm{A} _ {\lambda} ^ {*} \mathrm{A} _ {\lambda} ^ {\prime}).\tag{56}
\]

Since \(\hat{\tau}\circ\overline{{\mathcal{F}}}=\tau\) , formulas (53) and (55) imply that the mapping \(\overline{{\mathcal{F}}}\) is an isomorphism of Hilbert spaces from C{[}G{]} to F( \(\widehat{G}\) ).

The Schur orthogonality relations (VIII, p.~410) can be reformulated using Hilbertian inner products. Relations (21) and (24) then give the following assertions. For \(\lambda \in \widehat{\mathbf{G}}\) and \(x, x', y, y'\) in \(\mathrm{V}_{\lambda}\) , we have

\[
| \mathrm{G} | ^ {- 1} \sum_ {g \in \mathrm{G}} \overline {{\langle x | \pi_ {\lambda} (g) x ^ {\prime} \rangle}} _ {\lambda} \langle y | \pi_ {\lambda} (g) y ^ {\prime} \rangle_ {\lambda} = d _ {\lambda} ^ {- 1} \overline {{\langle x | y \rangle}} _ {\lambda} \langle x ^ {\prime} | y ^ {\prime} \rangle_ {\lambda}.\tag{57}
\]

If \(\lambda\) and \(\mu\) are two distinct elements of \(\widehat{\mathrm{G}}\) , then for \(x, x'\) in \(\mathrm{V}_{\lambda}\) and \(y, y'\) in \(\mathrm{V}_{\mu}\) , we have

\[
\sum_ {g \in \mathrm{G}} \overline {{\langle x | \pi_ {\lambda} (g) x ^ {\prime} \rangle_ {\lambda}}} \langle y | \pi_ {\mu} (g) y ^ {\prime} \rangle_ {\mu} = 0.\tag{58}
\]

For every \(\lambda \in \widehat{\mathbf{G}}\) , we choose an orthonormal basis \((e_{\lambda,i})_{1\leqslant i\leqslant d_{\lambda}}\) of \(\mathrm{V}_{\lambda}\) . For any \(g\in \mathbf{G}\) , we denote by \((\pi_{ij}^{\lambda}(g))\) the matrix of the endomorphism \(\pi_{\lambda}(g)\) of \(\mathrm{V}_{\lambda}\) with respect to this basis; we have

\[
\pi_ {i j} ^ {\lambda} (g) = \langle e _ {\lambda , i} | \pi_ {\lambda} (g) e _ {\lambda , j} \rangle_ {\lambda}.\tag{59}
\]

Since the endomorphism \(\pi_{\lambda}(g)\) is unitary, its inverse is equal to \(\pi_{\lambda}(g)^*\) , so that

\[
\overline {{\pi_ {i j} ^ {\lambda} (g)}} = \pi_ {j i} ^ {\lambda} (g ^ {- 1}).\tag{60}
\]

It then follows from formulas (22) of VIII, p.~409 and (25), p.~409 that the functions \((d_{\lambda})^{1 / 2}\pi_{ij}^{\lambda}\) , for \(\lambda \in \widehat{\mathbf{G}}\) , \(1\leqslant i\leqslant d_{\lambda}\) , \(1\leqslant j\leqslant d_{\lambda}\) , form an orthonormal basis of the Hilbert space \(\mathbf{C}[\mathrm{G}]\).*

